\documentclass[10pt,a4paper]{amsart}
\usepackage{amsmath,amssymb,amsthm,graphicx}
\usepackage{verbatim}
\newenvironment{reference}{\comment}{\endcomment}
%\newenvironment{reference}{}{}
\newenvironment{slogan}{\comment}{\endcomment}
\newenvironment{history}{\comment}{\endcomment}
% For commutative diagrams we use Xy-pic
\usepackage[all]{xy}

% We use 2cell for 2-commutative diagrams.
\xyoption{2cell}
% Theorem environments.
%
\theoremstyle{plain}
\newtheorem{theorem}[subsection]{Theorem}
\newtheorem{proposition}[subsection]{Proposition}
\newtheorem{lemma}[subsection]{Lemma}

\theoremstyle{definition}
\newtheorem{definition}[subsection]{Definition}
\newtheorem{example}[subsection]{Example}
\newtheorem{exercise}[subsection]{Exercise}
\newtheorem{situation}[subsection]{Situation}

\theoremstyle{remark}
\newtheorem{remark}[subsection]{Remark}
\newtheorem{remarks}[subsection]{Remarks}

\numberwithin{equation}{subsection}

% Macros
%
\def\lim{\mathop{\mathrm{lim}}\nolimits}
\def\colim{\mathop{\mathrm{colim}}\nolimits}
\def\Spec{\mathop{\mathrm{Spec}}}
\def\Hom{\mathop{\mathrm{Hom}}\nolimits}
\def\Ext{\mathop{\mathrm{Ext}}\nolimits}
\def\SheafHom{\mathop{\mathcal{H}\!\mathit{om}}\nolimits}
\def\SheafExt{\mathop{\mathcal{E}\!\mathit{xt}}\nolimits}
\def\Sch{\mathit{Sch}}
\def\Mor{\mathop{\mathrm{Mor}}\nolimits}
\def\Ob{\mathop{\mathrm{Ob}}\nolimits}
\def\Sh{\mathop{\mathit{Sh}}\nolimits}
\def\NL{\mathop{N\!L}\nolimits}
\def\CH{\mathop{\mathrm{CH}}\nolimits}
\def\proetale{{pro\text{-}\acute{e}tale}}
\def\etale{{\acute{e}tale}}
\def\QCoh{\mathit{QCoh}}
\def\Ker{\mathop{\mathrm{Ker}}}
\def\Im{\mathop{\mathrm{Im}}}
\def\Coker{\mathop{\mathrm{Coker}}}
\def\Coim{\mathop{\mathrm{Coim}}}

% Boxtimes
%
\DeclareMathSymbol{\boxtimes}{\mathbin}{AMSa}{"02}

%
% Macros for moduli stacks/spaces
%
\def\QCohstack{\mathcal{QC}\!\mathit{oh}}
\def\Cohstack{\mathcal{C}\!\mathit{oh}}
\def\Spacesstack{\mathcal{S}\!\mathit{paces}}
\def\Quotfunctor{\mathrm{Quot}}
\def\Hilbfunctor{\mathrm{Hilb}}
\def\Curvesstack{\mathcal{C}\!\mathit{urves}}
\def\Polarizedstack{\mathcal{P}\!\mathit{olarized}}
\def\Complexesstack{\mathcal{C}\!\mathit{omplexes}}
% \Pic is the operator that assigns to X its picard group, usage \Pic(X)
% \Picardstack_{X/B} denotes the Picard stack of X over B
% \Picardfunctor_{X/B} denotes the Picard functor of X over B
\def\Pic{\mathop{\mathrm{Pic}}\nolimits}
\def\Picardstack{\mathcal{P}\!\mathit{ic}}
\def\Picardfunctor{\mathrm{Pic}}
\def\Deformationcategory{\mathcal{D}\!\mathit{ef}}
\usepackage{cleveref}
\crefname{originalreference}{original reference}{original references}
\Crefname{originalreference}{Original reference}{Original references}
\makeatletter
\newlabel{situation-setup}{{02QL}{source}{}{sourcepointer-1}{}}
\newlabel{situation-setup@cref}{{[originalreference][0][]02QL}{[1][0][]source}{}{}{}}
\makeatother
\title{The bivariant group of a morphism between disjoint unions is the product of the bivariant groups of its pieces}
\author{The Stacks Project Authors}
\date{stacks-project:chow.tex:local-snapshot-20261012; GNU FDL 1.2 or later}
\begin{document}
\maketitle
\noindent\textit{Editorial note.} All delivered mathematical text is the author's own, copied from the pinned Stacks Project chapter file chow.tex; no statement or proof step was rewritten, reordered or removed.\par\medskip
\noindent\textit{Editorial note.} Preamble provenance: the pinned chapter file opens with \textbackslash{}input\{preamble\} and carries no preamble text of its own, so the project's single Stacks-wide preamble staged at sources/stacks-chow/preamble.tex is excerpted verbatim as three spans, lines 8-12 (the verbatim package together with the project's own reference, slogan and history display-environment definitions, including the author's commented-out older form of the reference environment), lines 14-18 (Xy-pic setup) and lines 154-217 (the native theorem declarations and mathematical macros). The lines 8-12 span is required because a counted statement or its surrounding text may use those display environments, which the project typesets as comments. Nothing was added, removed or redefined; the wrapper class and the amsmath/amssymb/amsthm/graphicx base packages are supplied by the assembler.\par\medskip
\noindent\textit{Editorial note.} Same-chapter pointers are located at their \textbackslash{}label lines inside the pinned chapter file; cross-chapter pointers are located in the project's own tag list staged in the snapshot (sources/stacks-chow/tags.txt), at the verbatim lines carrying their canonical tags.\par\medskip
\noindent\textit{Allowed standard prerequisites (editorial).}\begin{itemize}
\item The situation fixing the bivariant theory. Same chapter chow.tex, situation labelled situation-setup (statement, proof line 6624). Fixes the graded bivariant groups A\textasciicircum{}p(X -> Y) and the Chow groups on which they act; the rest of the proof uses no other citation.
\end{itemize}\medskip
\ifdefined\hypertarget\hypertarget{sourcepointer-1}{}\fi
\noindent\textit{Source reference.} Original reference 02QL: sources/stacks-chow/pinned-original-source.tex, source lines 1372--1372. Navigation pointer for \textbackslash{}ref\{situation-setup\} used inside the counted span (line 6624); the argument that consumes it is retained in the span.\par\medskip
\begin{lemma}
\label{lemma-disjoint-decomposition-bivariant}
Let $(S, \delta)$ be as in Situation \ref{situation-setup}.
Let $f : X \to Y$ be a morphism of schemes locally of finite type over $S$.
Assume we have disjoint union decompositions
$X = \coprod_{i \in I} X_i$ and $Y = \coprod_{j \in J} Y_j$
by open and closed subschemes
and a map $a : I \to J$ of sets such that $f(X_i) \subset Y_{a(i)}$.
Then
$$
A^p(X \to Y) = \prod\nolimits_{i \in I} A^p(X_i \to Y_{a(i)})
$$
\end{lemma}

\begin{proof}
Suppose given an element $(c_i) \in \prod_i A^p(X_i \to Y_{a(i)})$.
Then given $\beta \in \CH_k(Y)$ we can map this to the element of
$\CH_{k - p}(X)$ whose restriction to $X_i$ is $c_i \cap \beta|_{Y_{a(i)}}$.
This works because $\CH_{k - p}(X) = \prod_i \CH_{k - p}(X_i)$.
The same construction works after base change by any $Y' \to Y$
locally of finite type and we get $c \in A^p(X \to Y)$.
Thus we obtain a map $\Psi$ from the right hand side of the formula
to the left hand side of the formula.
Conversely, given $c \in A^p(X \to Y)$ and an element
$\beta_i \in \CH_k(Y_{a(i)})$ we can consider the element
$(c \cap (Y_{a(i)} \to Y)_*\beta_i)|_{X_i}$ in $\CH_{k - p}(X_i)$.
The same thing works after base change by any $Y' \to Y$
locally of finite type and we get $c_i \in A^p(X_i \to Y_{a(i)})$.
Thus we obtain a map $\Phi$ from the left hand
side of the formula to the right hand side of the formula.
It is immediate that $\Phi \circ \Psi = \text{id}$.
For the converse, suppose that $c \in A^p(X \to Y)$ and
$\beta \in \CH_k(Y)$. Say $\Phi(c) = (c_i)$. Let $j \in J$.
Because $c$ commutes with flat pullback we get
$$
(c \cap \beta)|_{\coprod_{a(i) = j} X_i} =
c \cap \beta|_{Y_j}
$$
Because $c$ commutes with proper pushforward we get
$$
(\coprod\nolimits_{a(i) = j} X_i \to X)_*
((c \cap \beta)|_{\coprod_{a(i) = j} X_i})
=
c \cap (Y_j \to Y)_*\beta|_{Y_j}
$$
The left hand side is the cycle on $X$ restricting to $(c \cap \beta)|_{X_i}$
on $X_i$ for $i \in I$ with $a(i) = j$ and $0$ else.
The right hand side is a cycle on $X$ whose restriction to $X_i$
is $c_i \cap \beta|_{Y_j}$ for $i \in I$ with $a(i) = j$.
Thus $c \cap \beta = \Psi((c_i))$ as desired.
\end{proof}

\end{document}
